% 本文件由 LaTeX 排版 Agent 根据有效 translation.json 自动生成。
% author=Ephraim the Syrian; work_id=3704; title=主显节赞美诗

\chapter{主显节赞美诗}

% structure: p0001 source=h1 target=omitted reason=page-title-already-rendered-by-parent

% structure: p0002 source=h2 target=section reason=relative-heading-rank; base=h2

\section{赞美诗 一}

应答。——\Emph{愿祢在祢的主显之日，从祢的羊群中得着赞美！}\par

1. 祂更新了诸天，因为愚人崇拜众光体：——\newline{}祂更新了大地，因为在亚当内它已荒废。——\newline{}祂所塑造的因祂的唾沫而成为新的：——\newline{}全能者用灵魂恢复了身体。\par

2. 你们再次聚集——羊群啊，不劳而获地领受洁净！——因为无需像厄里叟那样——在河中沐浴七次，也无需像司铎那样因洒水而疲惫。\par

3. 厄里叟七次洁净了自己，作为七神的一个奥迹；——牛膝草和血是强大的象征。——没有分裂的余地——\Emph{祂}不可与万有的主分离，祂是万有之主的圣子。\par

4. 梅瑟在\Place{玛辣}使苦水变甜——因为人民抱怨和发怨言：——\Emph{如此}他给了圣洗的一个预象——在其中生命的主使那些苦的变甜。\par

5. 云彩荫庇并挡住了营地的灼热——它显示了圣神的象征，祂在圣洗中荫庇你们——调和火焰，\Emph{使火焰不伤害}你们的身体。\par

6. 那时人民穿过大海，并显示了一个预象——就是你们在其中受洗的圣洗。人民穿过那里却不相信：——外邦人在这里受洗，相信了，并领受了圣神。\par

7. 圣言派遣了声音在祂来临之前宣告，——为祂预备祂来的道路——并聘娶新娘，直到祂来临——使她预备好，等祂来从水中迎娶她。\par

8. 预言的声音激动了不孕妇人的儿子——他出去在旷野中流浪并呼喊——\Quote{看！天国的儿子来了！——预备道路，使祂能进入并住在你们的居所！}\par

9. 若翰呼喊：\Quote{那在我以后来的，祂在我以前：——我是声音，但不是圣言；——我是火炬，但不是光——是在正义的太阳之前升起的星。}\par

10. 这个若翰在旷野中曾呼喊并说——\Quote{你们罪人，悔改你们的恶行——并献上悔改的果实；——因为看！那从莠子中簸扬麦子的祂来了。}\par

11. 光之赐予者得胜了，并标记了一个奥迹，藉着他上升的度数：——看！他上升后已有十二天，——而今天这是第十三天：——祂，圣子，和祂十二位的一个完美奥迹！\par

12. 黑暗被克服，为显明撒殚被克服——而光明得胜，使他宣告——长子得胜：黑暗被克服——与黑暗之灵一同，而我们的光明与光之赐予者一起得胜。\par

13. 在高处和深渊，圣子有两个先驱。——光之星从高处宣告祂——若翰同样从低处宣讲祂：——两个先驱，属地的和属天的。\par

14. 光之星，违反自然，突然照耀——小于太阳却大于太阳。——它在显明的光中小于他——而在隐秘的权能中大于他，因为它的奥迹。\par

15. 光之星在那些在黑暗中的人中放射光芒——并引导他们如同他们是瞎子——使他们前来遇见大光：——他们奉献了礼物，领受了生命，朝拜了，并离去了。\par

16. 从高处来的\Emph{先驱}显示祂的本性来自至高者；——同样，从低处来的\Emph{那位}显示祂的身体来自人类，大奇事！——使祂的天主性和祂的人性藉着他们被宣告！\par

17. 因此，谁若认为祂是属地的，光之星——将说服他祂是属天的：谁若认为祂是属灵的，——这个若翰将说服他祂也是有身体的。\par

18. 若翰与他的父母靠近并朝拜了太阳——荣光停留在祂的面上。——他没有像在母腹中那样激动。——大奇事！他在这里朝拜，而在那里欢跃！\par

19. 整个受造物为祂成了一个口，并关于祂呼喊。——贤士在他们的礼物中呼喊；——不孕者与他们的孩子呼喊——光之星，看！它在空中呼喊：\Quote{看，君王的儿子！}\par

20. 诸天开了，水涌出，鸽子在荣耀中！——圣父的声音比雷更强——当它说出这话：\Quote{这是\Emph{我的爱子}}；——守望者带来了消息，孩子们在他们的贺撒纳中欢呼\Emph{祂}。\par

% structure: p0024 source=h2 target=section reason=relative-heading-rank; base=h2

\section{赞美诗 二}

（与赞美诗十三\WorkTitle{论耶稣圣诞}几乎相同。）\par

（应答。——\Emph{愿祢受赞美，祢在这庆节中使万众欢腾！}）\par

1. 在那被称为\OriginalTerm{Semha}{Semha}的君王时代——我们的主在希伯来人中被显明。——因此\OriginalTerm{Semha}{Semha}和\OriginalTerm{Denha}{Denha}一同为王——地上的君王和高处的圣子——愿\Emph{祂的权能受赞美}！\par

2. 在那登记人民以征税的君王的日子里，——我们的救主降下，并在生命册上登记了人民；祂写了，也被写了——在高处祂写了我们，在地上祂被写；光荣归于祂的圣名！\par

3. 祂的诞生是在那名叫\OriginalTerm{Semha}{Semha}的君王的日子里。——预象与真理彼此相遇——君王与君王，\OriginalTerm{Semha}{Semha}与\OriginalTerm{Denha}{Denha}。——那个王国背负了祂的十字架；愿那背起它的祂\Emph{受赞美}！\par

4. 三十年祂在地上贫困地居住。——赞美的声音用各种韵律——让我们为我们的主的岁月编织，我弟兄们——三十个冠冕为三十年；愿祂的数字\Emph{受赞美}！\par

5. 在第一年，宝藏之母，充满祝福，——让革鲁宾与我们一同感恩，他们背负——在荣耀中的圣子，祂放弃了祂荣耀的状态——并劳作，找到了那迷失的羊——愿感恩\Emph{归于祂}！\par

6. 在第二年，让色辣芬与我们一同倍增感恩——他们曾向圣子呼喊\Quote{圣}，却转身看见祂——在不相信的人中受辱。——祂忍受了嘲笑，并教导\Emph{我们}光荣；愿光荣归于祂！\par

7. 在第三年，让\Person{弥额尔}和他的众军与我们一同感恩——他们惯于在高处服侍圣子——\Emph{却}看见祂在地上服侍。——祂洗了\Emph{人的}脚，洁净了\Emph{人的}灵魂；愿\Emph{祂的温良受赞美}！\par

8. 在第四年，让诸天与我们一同感恩！对圣子而言\Emph{太}狭窄，它将要爆裂，看着——祂躺在被轻视的\Person{匝凯}的床上。——祂充满了床，也曾充满了诸天；——愿感恩\Emph{归于祂}！\par

9. 在第五年，愿那以热力灼烧大地的太阳，感谢我们的太阳，因祂收敛了自己的广阔，并调和了自己的力量，使\Emph{眼睛}得以忍受看见祂——那纯洁灵魂的内在眼睛。愿祂的光辉\Emph{受}赞颂！\par

10. 在第六年，愿所有空气同我们一道感谢——在其广阔之中万有欢跃。它看见自己的伟大之主成为了——一个卑微怀抱中的小婴儿。愿祂的尊荣\Emph{受}赞颂！\par

11. 在第七年，愿云彩和风与我们同吹\Emph{号角}——它们的露珠洒在花朵的面庞——它们却看见圣子收敛了自己的光明——忍受了讥讽和可耻的唾弃——愿祂的救恩\Emph{受}赞颂！\par

12. 又在第八年，愿受造界献上光荣——从它的泉源中，果实汲取滋养。它朝拜了，当它看见圣子在怀抱中——纯洁的\Emph{婴儿}由纯洁的乳汁养育。愿祂的悦纳\Emph{受}赞颂！\par

13. 在第九年，愿大地献上光荣——它的怀中得灌溉，便生出根。——它看见玛利亚是一片未受浇灌的土壤——她所结的果实是一片汪洋大海。愿向祂\Emph{献上}欢呼！\par

答：愿光荣归于祢，万有之主之子，祢赐生命给万有！\par

14. 在第十年，愿西乃山献上光荣——它在主面前曾熔化！它看见有人拿起石头攻击它的主：但祂却拿起石头——把教会建造在磐石之上；愿祂的建筑\Emph{受}赞颂！\par

15. 在第十一年，愿大海感谢——圣子那测量过它的手！它惊奇地看见祂降下，在卑微的水中被洗——祂，那洁净受造界的。愿祂的胜利\Emph{受}赞颂！\par

16. 在第十二年，愿圣殿感谢——它看见那孩童坐在——长老中间：经师们静默不语——如同逾越节的羔羊在节期中哞叫。愿祂的赎罪\Emph{受}赞颂！\par

17. 在第十三年，愿冠冕与我们一同感谢——那位得胜并被荆棘冠冕加冕的君王：祂为人编织了一顶强大的冠冕，在祂的右边。愿那派遣祂的\Emph{受}赞颂！\par

18. 在第十四年，愿埃及的逾越节感谢——那来了为众人成就逾越的逾越节——它代替法老淹没了军团——代替骑马者淹没了魔鬼。愿祂的报应\Emph{受}赞颂！\par

19. 在第十五年，愿羊群中的羔羊感谢——我们的主没有像\Emph{梅瑟那样}宰杀它——而是以祂的血救赎了人类。——祂是万有的牧者，为众人而死；愿那生祂的\Emph{受}赞颂！\par

20. 在第十六年，愿奥秘中的谷粒感谢——那位农夫，祂把自己的身体当作种子——种在腐蚀万物的贫瘠土壤中。它竟变得肥沃，结出了新的食粮；愿那位纯洁者\Emph{受}赞颂！\par

21. 在第十七年，愿葡萄树感谢我们的主——真理的葡萄园，\Emph{其中}灵魂如同枝条。祂赐平安给这葡萄园，却荒废了那结出野葡萄的葡萄园；愿那拔除者\Emph{受}赞颂！\par

22. 在第十八年，愿我们的酵母感谢——那真理的酵母，它渗透并吸引——所有心思，使它们\Emph{成为}——同一心思，同一道理。愿祂的道理\Emph{受}赞颂！\par

23. 在第十九年，愿盐为祢的身体感谢。——噢，有福的婴孩，是灵魂——作为身体的盐，而信德——是灵魂的盐，借此它得以保存；愿祢的保存\Emph{受}赞颂！\par

答：光荣归于祢的显现，噢，天主而人！\par

24. 在第二十年，愿世间的财富与我们一同感谢——那些\Emph{成全}的人已经丢弃并放弃了它——因为那\Quote{祸哉}；而去爱慕——贫穷，因为它的真福。愿那渴慕它的\Emph{受}赞颂！\par

25. 在第二十一年，愿那因圣子的奥秘而变甜的水感谢。在参孙的蜜中，万民尝到了里面的苦味，那苦味毁灭了他们：他们在救赎他们的十字架中得到了生命；愿它的甘甜受赞颂！\par

26. 在第二十二年，愿武器和刀剑感谢——因为它们不能杀死我们的仇敌。——祢就是那杀死他的，正如祢就是那复原——西满的剑削掉的耳朵的；愿祢的医治\Emph{受}赞颂！\par

27. 在第二十三年，愿驴同样感谢——它献出了小驴驹给祂骑乘——祂也同样开了野驴的口——后代向\Emph{祂}献上赞美；愿祢的赞美\Emph{受}赞颂！\par

28. 在第二十四年，愿财富感谢圣子！——宝藏们对宝藏之主感到惊讶——祂如何在穷人中间成长。——祂使自己贫穷，为能使众人富足；愿祂的分享\Emph{受}赞颂！\par

29. 在第二十五年，愿依撒格感谢圣子——祂在山上救他免于刀下——并代替他成为那将被宰杀的羔羊。——那必朽的逃脱了，而祂，那赐生命给万有的，却死了；愿祂的奉献\Emph{受}赞颂！\par

30. 在第二十六年，愿梅瑟与我们一同感谢——他害怕并逃离了杀人者；——愿他感谢圣子，因为正是祂亲身——进入阴府，劫掠了它并出来；愿祂的复活\Emph{受}赞颂！\par

31. [在第二十七年，愿雄辩的演说家感谢圣子，因为他们找不到——使我们能在判决中得胜的方法：——祂在审判中沉默，却使我们的判决得胜；愿向祂\Emph{献上}掌声！]\par

32. 在第二十七年，愿所有审判者感谢——他们作为义人处死了作恶者；——愿他们感谢圣子，祂代替恶人——作为善人而死，尽管祂是义者之子；愿祂的仁慈\Emph{受}赞颂！\par

33. 在第二十八年，愿他们感谢圣子——所有那些不曾从掳掠者手中拯救我们的勇士。有一位当受朝拜——祂被杀害、被擒拿，却拯救了我们；愿祂的拯救\Emph{受}赞颂！\par

34. 在第二十九年，愿约伯与我们一同感谢——他为主担当了苦难：——但我们的主为我们担当了——唾弃、鞭打、荆棘和钉子；愿祂的怜悯受赞颂！\par

35. 在那第三十年，愿他们与我们一同感谢——那些因祂的死亡而活了的死者——那些因祂的十字架而悔改的活人——以及那在祂内得以和好的高天和深渊！愿祂和祂的圣父\Emph{受}赞颂！\par

% structure: p0064 source=h2 target=section reason=relative-heading-rank; base=h2

\section{第三首赞美诗}

(答唱：— \Emph{基督以圣化圣油，看哪！祂正在印封祂羊群中新生的羔羊！})\par

1. 基督与圣化圣油相连；隐秘与可显者相混合：圣化圣油有形地傅抹——基督秘密地印封，那些新生、属灵的羔羊，祂双重胜利的奖品；因为祂由圣化圣油生了它，又由水生了它。\par

2. 你们的圣秩何其尊高！因为那曾为罪妇者，如婢女般，傅抹了主的脚。\Emph{但}为你们，基督却好似\Emph{祂的}仆人，藉祂仆人之手，印封并傅油你们的身体。祂身为群羊之主，亲自印封祂的羊群是相宜的。\par

3. 既然那罪妇需要宽恕，傅油对她便是献礼，藉此她的爱使她的主和好。但你们身为羊群，处于亵渎和不奉信者中，真理藉圣化圣油成为你们的印号，使你们与迷途者分离。\par

4. 祂藉割损的昔日印号，将百姓从万民中分别出来；但藉傅油的印号，祂将万民从百姓中分别。当万民误入歧途时，祂将百姓从万民中分别；如今当百姓背离祂时，祂便将万民从那里分别。\par

5. \Person{纳阿曼}取走净土的尘土，回到本地；好藉这圣土，与不洁者分别并认出。基督的圣化圣油将奥迹之子与陌生人分开：藉此，内里者被分别，与外在者得以辨认。\par

6. \Person{厄里亚}曾增多的油，可用口品尝；因为那罐是寡妇的，并非圣化圣油的油罐。我主在罐中的油，不是口腹的食物：那曾如外狼的罪人，它使之成为羊群中的羔羊。\par

7. 温良谦逊者的圣化圣油，改变倔强者，使之相似其主。外邦人原是豺狼，畏惧\Person{梅瑟}严厉的棍杖。看哪！圣化圣油印封他们，将豺狼变成羊群！那曾逃避棍杖的豺狼，看哪！他们如今投奔了十字架！\par

8. 橄榄叶到达，作为傅油的预象带来；方舟之子欢喜迎接它，因它带来救援的福音。同样你们也欢喜迎接这神圣傅油。罪人的身体因它喜乐，因它带来救援的福音。\par

9. \Person{雅各伯}曾傅油于石上，当他立石为记号时，使这石成为他与天主之间的凭证，并在那里献上什一；看哪！这正是你们身体的预象，\Emph{如何}因圣化圣油而被印封\Emph{为}圣，成为天主的圣殿，在那里祂将因你们的祭献受事奉。\par

10. 当\Person{梅瑟}印封并傅油了\Person{亚郎}司祭的儿子们，烈火吞噬了他们的身体；烈火却保全了他们的祭衣。但你们，我的弟兄，你们有福了，因为恩宠之火已降临，彻底焚尽了你们的过犯，洁净并圣化了你们的身体！\par

11. 至于\Person{亚郎}的傅油，我的弟兄们，那是牲畜卑贱的血，洒在祭台的角上。真理的傅油却是这样：其中活而全赐生命的血，内在地洒在你们身上，混合在你们的理智中，渗透在你们的内心深处。\par

12. 受傅的司祭曾献上被宰杀的牲畜身体；你们，受傅且超凡者，你们的祭品\Emph{是}你们\Emph{自己的}身体。受傅的肋未人曾献上牲畜的内脏：你们已超越肋未人，因为你们献上了自己的心。\par

13. 百姓的傅油——是基督的预影；他们的棍杖是十字架的奥迹；他们的羔羊是独生子的预象；他们的帐幕是你们教会的奥迹；他们的割损是你们印号的标记。在你们美好之事物的荫下，古时的百姓曾安坐。\par

14. 真理恰如一棵大荫树：曾将荫影投在百姓上；将根扎在万民中。百姓住在它荫下，它的荫影即是其奥迹；但外邦人寄居在它的枝上，摘食它的果实。\par

15. \Emph{至于}\Person{撒乌耳}受傅为王，其香气越甘甜，他心肠的臭气越甚。圣神击打他，随即离去。你们所领受的傅油更伟大；因为你们的理智是香炉，在你们的殿宇中圣神欢跃，你们将永远成为祂的内室。\par

16. \Emph{至于}\Person{达味}的傅油，我的弟兄们，圣神降临，在祂所喜悦的人心中发出甘香；他内心的香气正如他行为的香气。圣神住在他内，在他内歌唱。你们所领受的傅油更伟大，因为圣父、圣子及圣神，已发动并降临，住在你们内。\par

17. 古时癞病人得洁净，司祭常用油印封他，并领他到水泉。预象已过，真理已来；看哪！你们已用圣化圣油得印封，在圣洗中得成全，在羊群中得混合，由身体得滋养。\par

18. 有哪个癞病人得了洁净，还转回贪恋他的癞病？你们已脱去过犯——要抛弃它！无人再穿上他已脱去的癞病。它已落下沉没——不可再拉出来！它已消耗磨损——不可再更新！愿腐朽不再临于你们，你们是基督的圣化圣油所傅抹的人！\par

19. 泥土塑造的器皿，得水之美，受火之力；\Emph{但}若失手，便毁坏了，不能再重新制作。你们是恩宠的器皿；要提防它，甚至要提防公义，因为公义不赐两次更新。\par

20. 你们何等相似那被鱼吐出的先知！吞灭者已交还你们，因他受制于那控制鱼的大能。\Person{约纳}为你们\Emph{是}一面镜子，既然鱼不再吞他，愿吞灭者也不再吞你们：既已吐出，当如\Person{约纳}！\par

21. \Person{玛利亚}曾将珍贵的香膏倒在我主的头上；它的芬芳充满全屋。同样你们傅油的香气，已芬芳并馨香诸天，达于高天的守望者；使撒殚不悦，它的香气\Emph{是}压倒性的；对天主，它的气味是甘甜的。\par

22. 旷野中的群众如同没有牧人的羊群。仁慈者成了他们的牧人，为他们倍增了饼的牧场。是的，你们这些成全的人有福了，你们被盖上了基督羔羊的印号，你们配得上祂的身体和鲜血；牧人自己成了你们的牧场！\par

23. 祂用水变成了酒，在筵席上赐给青年们饮用。对于守斋的你们，傅油比饮酒更美。在祂的酒中，订婚者结合；在祂的油中，已婚者得以圣化。祂的酒带来结合，祂的油带来圣化。\par

24. 基督的羊群欢跃，为接受生命的印号，这君王的徽标永远驱散罪恶。恶者因您的徽标溃败，不法因您的记号四散。来吧，群羊，接受你们的印号，它驱散那吞吃你们的！\par

25. 来吧，羔羊们，接受你们的印号，因为真理就是你们的印号！这印号分别家中的人与外人。刀兵同样割损了悖逆者和哈加尔的儿子。如果割损是羊的记号，看！山羊也被盖上了记号。\par

26. 但是你们，新生的羊群，已经脱去了豺狼的行径，如同羔羊一样相似于羔羊。一位以改变改变了万有；羔羊将自己交与豺狼，任其被宰杀；豺狼冲来吞吃了祂，却变成了羔羊；因为牧人变成了羔羊，同样，豺狼忘记了其本性。\par

27. 也求祢以仁慈眷顾我！不要让那山羊、左道之子的印号烙在我身上！不要使祢的羊变成山羊！因为我虽不足以成义，也不愿成为罪人。我主，求祢转目不看我所作的，也不只追究我所愿的。\par

28. 愿从书写者和宣讲者，从听者和受印者，荣耀归于基督，并藉着祂归于圣父，受赞颂！祂赐言语给说话者，赐声音给宣讲者，赐理解给听者，并祝圣傅油给受印者。\par

% structure: p0094 source=h2 target=section reason=relative-heading-rank; base=h2

\section{赞美诗 四}

（应—— \Emph{在水中涂抹无量过犯的那一位是可赞美的！}）\par

1. 降下吧，我受印的弟兄们，穿上我们的主，—— 并重新加入祂的族系，因为祂是伟大族系之子，—— 正如祂在祂的圣言中所说的。\par

2. 祂的本性来自高处，祂的外衣来自下方。—— 凡脱去自己外衣的，那外衣便与祂的外衣永远相融。\par

3. 你们也在水中，从祂接受那外衣——它不磨损也不失落，因为那是穿上它的人永远穿着的衣裳。\par

4. 但蒙福的司铎是两者之间的中保：—— 盟约将在祂面前订立，祂是祂主的中保，—— 也是我们这一方的保证。\par

5. 神性在水中，看哪！祂已掺入了祂的酵母——为尘土所造之物，那酵母高举他们——神性与他们结合。\par

6. 因为那是主的酵母，能渗入奴仆中——并高举他得自由；它使奴仆加入祂的族系——万有之主的族系。\par

7. 因为那在水里穿上祂、使万有自由的奴仆——虽在地上是奴仆，却是高天自由人的儿子——因为他穿上了自由。\par

8. 那在水里穿上那天使的自由人——成了仆役的同僚，为使他相似那成为众仆之仆的主。\par

9. 那使万有富足者降下，穿上了贫穷——为将水中宝库所藏的丰富分给穷人。\par

10. 那在下者又穿上了在水里赐予一切伟大的那一位——虽在愚者眼中卑微，在守望者眼中却伟大——因他穿上了伟大。\par

11. 因为正如那位伟大者，因爱而成为卑微——被不信者迫害，被守望者崇拜——祂成为卑微，也使卑微者伟大。\par

12. \Emph{这样}，让伟大者成为卑微，好使卑微者在他内伟大：——让我们相似那超越万有却成为小于万有的那一位：——祂成为卑微，并使众人伟大。\par

13. 那在水里穿上伟大者的温良人——虽面容谦和，其辨别力却极大——因为那超越万有者住在他内。\par

14. 因为谁能找到轻看荆棘丛的人呢？——那被轻看而谦卑的荆棘，其中住着在火中的威严。\par

15. 又有谁能找到轻看梅瑟的人呢？——那温良而拙口的，当那超越的荣耀——停驻在他的温良之上时？\par

16. 轻看他的，就是轻看他的主；那轻看他的恶人——大地发怒吞没了他们；那藐视祂的肋未人——烈火忿怒地吞噬了他们。\par

17. 基督曾命令他们：\Quote{你不可叫他\InnerQuote{拉加}}，那受洗并穿上祂的人；因为凡轻看那被轻看的，就是与那大能者一同轻看。\par

18. 在伊甸园和世界中，都有我们主的比喻；——有什么言语能收集祂奥秘的比拟呢？——因为祂完全在万物中被描绘。\par

19. 祂被记载于圣经中；印刻在自然上——祂的冠冕在君王身上描绘，祂的真理在先知身上，祂的赎罪在司祭身上。\par

20. 祂在梅瑟的棍杖中，在亚郎的牛膝草中——在达味的冠冕中：先知拥有祂的相似，宗徒拥有祂的福音。\par

21. 启示看见了祢，箴言盼望了祢——奥秘等候了祢，比拟问候了祢，比喻显示了祢的预像。\par

22. 梅瑟的盟约期盼那福音：——古时的一切都飞翔并落在新盟约上。\par

23. 看哪！先知们已将他们的荣耀奥秘倾倒在祂身上——司祭和君王已将他们的奇妙预像倾倒在祂身上：——他们全都将\Emph{它们}倾倒在祂全身上。\par

24. 基督以祂的教导胜过并超越了奥秘——以祂的解释胜过了比喻；如同大海从中——容纳一切江河。\par

25. 因为基督就是大海，祂能容纳——那从圣经中涌出的泉源和溪流、江河与川流。\par

% structure: p0121 source=h2 target=section reason=relative-heading-rank; base=h2

\section{赞美诗 五}

（应—— \Emph{那为了亚当子孙的赎罪而设立圣洗的那一位是可赞美的！}）\par

1. 降下吧，我的弟兄们，从圣洗的水中穿上圣神——与那服侍神性的诸灵结合！\par

2. 因为看哪！祂是那隐秘的火，也秘密地以三个属灵的名字封印祂的羊群，恶者在此被驱逐。\par

3. \Person{若翰}呼喊说：\BibleQuote{这是天主的羔羊}，——由此表明外邦人是\Person{亚巴郎}的子孙。\par

4. 这就是那为我们的救主作证的一位：祂要以火和圣神施洗。——看哪！火和圣神，我的弟兄们，在真理的圣洗中。\par

5. 因为圣洗比\Place{约旦河}这条小河更大——因为在众水与油之流中，所有人的恶行被洗除。\par

6. \Person{厄里叟}藉着七次\Emph{洗涤}，洁净了\Person{纳阿曼}的癞病；——在圣洗中，灵魂里隐秘的恶行被洁净。\par

7. \Person{梅瑟}在海中为百姓施洗，却未能——洗淨他们内心的邪恶，那心充满了恶行的污秽。\par

8. 看哪！司铎以\Person{梅瑟}的样式，涤除灵魂的污秽——并以傅油之油，看哪！他为新生的羔羊盖上印记，为天国。\par

9. \Person{撒慕尔}在百姓中傅油\Person{达味}为王：——但看哪！司铎傅油你们，为在天国中作嗣子。\par

10. 因为\Person{达味}所穿戴的铠甲，在受傅之后，他作战——并击倒了那试图征服以色列的巨人。\par

11. 看哪！在基督的\OriginalTerm{圣油}{chrism}中，以及从水中来的铠甲中——那试图征服外邦人的恶者的骄傲被谦抑了。\par

12. 从磐石流出的水，解了百姓的渴。看哪！在基督的泉源中，万民的渴被解了。\par

13. \Person{梅瑟}的杖开启了磐石，溪流涌出；那些因口渴而衰竭的人，因这饮料而复苏。\par

14. 看哪！从基督的肋旁流出了赋予生命的河流。——疲倦的外邦人饮了，并在其中忘记了他们的痛苦。\par

15. 愿祢的露珠洒在我的卑贱上，我的罪愆在祢的血中得到补赎！——我将在祢的右边，我的主，与祢的诸圣相联！\par

% structure: p0138 source=h2 target=section reason=relative-heading-rank; base=h2

\section{赞美诗 6}

(应和。——\Emph{愿那位受洗以使你受洗，好使你被赦免罪过者，受赞颂。})\par

1. 圣神从高天降下——藉着祂的覆翼祝圣了众水。——在\Person{若翰}的洗礼中——\Emph{祂}越过其余，只住在一人身上：——但现在祂降下并住在——所有由水而生的人身上。\par

2. 在\Person{若翰}所施洗的众人中——圣神只住在一位身上：——但现在祂飞翔降临——为住在许多人身上——并且当一个个上来时——祂爱他并住在他身上。\par

3. \Emph{这}是超越一切的奇事！——祂下到水中并受了洗。——众海都称它为有福的——祢在其间受洗的那条河：——甚至天上的水也嫉妒，——因为它们不配作祢的浴盆。\par

4. \Emph{这}是奇事，我的主，如今也是——当泉源充满水时——只有圣洗的水——能够赎罪。——海中的水是强大的——然而对于赎罪，它太弱了。\par

5. 我的主，祢的大能若住在——卑微者内，便高举他——如同王权若住在——旷野内，便赐予平安。——水因祢的大能战胜了——罪恶，因为生命环绕了它。\par

6. 羊群看见——手伸近为她们施洗，便欢跃。——领受吧，羊群，你们的印记；进入并汇入羊群！——因为今天\OriginalTerm{守望者}{Watchers}为你们欢欣，——胜过为整个羊群。\par

7. 天使与守望者欢欣——由于那由圣神和水所生的：——他们欢喜，因为\Emph{藉着}火与\Emph{藉着}圣神——属血肉的变成了属灵的。——\OriginalTerm{色辣芬}{Seraphins}颂唱\BibleQuote{圣}，也因成圣的人增多而欢欣。\par

8. 因为看哪！天使因一个悔改的罪人而喜乐：——何况现在他们喜乐——因为在所有教会和会众中——看哪！圣洗正将——属天的从属地中诞生出来！\par

9. 受洗者上来时被祝圣；——受印者下去时被赦免。——上来的人已穿上了光荣——下去的人已脱去了罪恶。——\Person{亚当}在一瞬间失去了他的光荣——你在瞬间被穿上了光荣。\par

10. 一所尘土造的房屋倒塌后——能藉水得以更新：——\Person{亚当}的尘土之身体——因水而更新了。——看哪！司铎们如同建筑者——重新更新你们的身体。\par

11. 羊毛的这一个大奇事——就是它能接受任何染料，——如同心灵\Emph{接受}任何言论。——它以其染色的名称被称呼——正如你们这些曾受洗的\Quote{听众}，——获得了\Quote{领受者}的名称。\par

12. \Person{厄里叟}藉着那隐秘的名——\Emph{甚至}祝圣了普通的水。——癞病人在其中公开洗浴——并藉着那隐秘的大能洁净了：——癞病在水中被消除，如同过犯在圣洗中一样。\par

13. 今天，看哪！你们的过犯被涂抹——你们的名字被记录。——司铎在水中涂抹——基督在天上记录。——藉着涂抹与记录——看哪！你们的喜乐加倍了。\par

14. 看哪！仁慈今天破晓——从边界延伸到边界：——太阳沉落，仁慈破晓。——公义收回了她的愤怒；恩宠扩展了她的爱——看哪！她白白地赦免并赋予生命。\par

15. 先前在\Emph{羊栈}中的羊——看哪！她们急速出来迎接——\Emph{于其中}新添加的小羊。——\Emph{他们是}洁白\Emph{并}穿着白衣——内外洁白\Emph{是}你们的身体，如同你们的衣裳。\par

16. 从每张口：\Quote{你们\Emph{是有福的},}——从四方：\Quote{你们\Emph{是有福的}.}——罪恶从你们身上被驱逐——圣神住在你们身上。——恶者面容忧伤——良善的天主使你们面容喜乐。\par

17. 你们白白领受的恩赐——要不断守护它：——这珍珠若丢失——不能再寻回——因为它如同童贞——一旦失落，便无处可寻。\par

18. 愿你们远离一切污秽——因你们白衣的力量而受保守！——愿那自由被自己玷污的人——能以他的眼泪洗淨自己！——至于我，是团体的仆人——愿团体的祈祷为我赢得赦免！\par

19. 致在言辞上劳苦的作者——愿安息中的和好归于他！——致以声音劳苦的教师——愿藉恩宠获得的宽恕归于他！——致在施洗中劳苦的司铎——愿义德的冠冕临于他！\par

20. 从每一个口中，一心一意——从下界和上界——\Person{守视者}、\Person{革鲁宾}和\Person{色辣芬}——受洗者、受印者和听者——让我们各人高声呼喊说：\Quote{光荣归于我们庆节之主！}\par

% structure: p0160 source=h2 target=section reason=relative-heading-rank; base=h2

\section{第七首赞歌}

（应答。—— \Emph{祂赎清了你们的罪，使你们配得领受祂的圣体！}）\par

1. \Person{雅各伯}的羊群下来——围绕在水井旁。——在水中，它们披上了那被\Emph{它}覆盖的木杖的相似。——这些是奥秘和十字架的预象，——比喻在其中得到解释。\par

2. 在\Emph{这些}木杖中显出了象征——在羊群中，显出了比喻。——十字架在木杖中被预表，在羊群中则是\Emph{人的}灵魂。——他的木杖是我们之木的奥秘；——同样，他的羊群是我们羊群的奥秘。\par

3. 基督的羊群欢乐——围绕在圣洗池旁——在水中，他们披上了——生活而美丽的十字架的肖像——万有都凝视其上，——而整个肖像印在他们众人身上。\par

4. 在井旁，\Person{黎贝加}接受了——耳环和手上的首饰。——基督的新娘已穿上了——来自水中的珍宝：——她手上是生活的身体，——耳中是许诺。\par

5. \Person{梅瑟}打了\Emph{水}，饮了——\Person{耶特洛}——罪孽之司祭的羊群。——但我们的牧者已为自己的羊群付洗——祂是真理的大司祭。——在井旁羊群是哑的——但在这里，羊有言语。\par

6. 子民经过水而受了洗：——子民上了干地，却成了\Emph{如同}外邦人。——诫命在他们耳中无味——玛纳在他们器皿中腐坏。——请吃这生活的身体，——这赋生命予万有的生命之药！\par

7. \Person{梅瑟}对\Person{罗特}的子孙说：\Quote{给我们水喝，我们付钱——只让我们从你的边界经过。}——他们拒绝了道路，也拒绝了暂时的水。——看！活水白白\Emph{赐予}——以及通往\Place{伊甸园}的道路！\par

8. \Person{基德红}从水中为自己拣选了——在战斗中得胜的人。——你已下到得胜的水中：——上来并在战斗中得胜！——从水中领受和好，——从战斗领受冠冕！\par

9. 你们受了洗，领受你们的灯——如同\Person{基德红}家的灯；——用你们的灯征服黑暗，——用你们的贺三纳征服沉默！——同样，\Person{基德红}在战斗中——因呐喊和火焰而获胜。\par

10. \Person{达味}王曾渴望——井里的水，他们给他取来了——但他没有喝，因为他看出那是用人的血买来的。——你在水中畅饮——那用天主的血买来的水。\par

11. 先知从\Place{厄东}看见——天主\Emph{来临}，如同\Emph{压榨葡萄者}。——祂预备了义怒的酒醡——祂践踏了万民，拯救了子民。——祂已转而设立了圣洗；——万民存活，子民却归于无有。\par

12. \Person{耶肋米亚}在河中埋藏了——那被损坏的麻布腰带——而[子民]变老并朽坏。——……——那些朽坏而损坏的万民——因水而披上了新。\par

13. 在\Place{史罗亚}，那蒙福的溪流中——司祭们为\Person{撒罗满}傅了油。——他的青年曾受尊荣——他的老年却被鄙视。——藉着纯洁的水，你们已披上了——天堂的纯洁。\par

14. 那从露水中干燥的羊毛——\Place{耶路撒冷}在其中被预表：——那装满水的盆——圣洗在其中被预表。——那羊毛按预像的方式是干的；——这盆按象征的方式是满的。\par

15. 疲倦的身体在水中——洗浴并从劳苦中恢复精力。——看！那隐藏着——恢复、生命和喜悦的盆。——在其中疲倦的\Person{亚当}得到了安息——他把劳苦带进了受造界。\par

16. 身体中汗的泉源——被设定以\Emph{保护}免于发烧：——圣洗的泉源——被设定以\Emph{保护}免于火焰。——这是那能——熄灭\Place{地狱}的水。\par

17. 那穿越旷野的人——以水作为盔甲——对抗那征服一切的干渴。——请下到基督的泉源——在你们的肢体中领受生命——作为对抗死亡的盔甲。\par

18. 同样，潜水者从海中——取上珍珠。——请受洗并从水中取上——那隐藏其中的纯洁——作为宝石镶在天主性冠冕上的珍珠。\par

19. 水手在器皿中——贮存甜水作为储备——在海洋中他贮存并保存它，在苦水中保持甜水。——同样，在罪的洪流中——请保存圣洗的水。\par

20. \Person{撒玛黎雅妇人}对我们的主说：\Quote{看！的确，这井很深。}——圣洗\Emph{虽高}——却因仁慈俯就了我们：——因为和好来自天上——降临到罪人身上。\par

21. \Quote{谁喝了我所赐给他的水，他将永远不再渴。}——为此，我亲爱的，你们要渴慕这神圣的圣洗——你们将永不再渴——以至于你们无需再领受另一个圣洗。\par

22. 在\Place{史罗亚}的洗礼中——瞎子洗涤了，他的眼球——被水开启并照亮——他脱去了那\Emph{是}在它们上的黑暗。——你们脱去了隐藏的黑暗——从水中穿上了光明。\par

23. \Person{比拉多}洗了手——以免列入杀害者之中。——你们沐浴了身体——手连同口。——请进去，列入那吃的人中——因为这生命之药赋生命予万有。\par

24. \Quote{来跟随我，我要使你们成为渔人的渔夫。}——因为代替那易朽的渔获——他们捕捞了那\Emph{是}永远的渔获。——他们曾捕鱼为死——如今却给人付洗，赐生命予那将要死的人。\par

25. 一百五十条鱼——被\Person{西满}的网从水中捕得——但藉着他的宣讲，从圣洗的怀抱中——捕得了成千上万的人——一网天国之子。\par

26. 看！我们的司铎如同渔夫——站在浅水之上——他从中捕得了大网——各种形状，各种类别——他拉起渔获，为使\Emph{它}靠近——至高的万王之王。\par

27. \Person{西满}取了鱼，拉上来——它们被带到主前：——我们的司铎从水中取出了——藉着他从\Person{西满}领受的手，——贞女和洁身男子被带到——诸节之主的庆节中。\par

28. 我用祢的仁慈恳求祢宽恕我——因为祢也曾以仁慈起誓——\OriginalTerm{辣步尼}{Rabboni}，\Quote{在死者的死亡中——我不喜欢，却喜欢他的生命。}——祢已起誓，我也恳求：——啊，起誓的祢，宽恕那恳求者吧！\par

% structure: p0190 source=h2 target=section reason=relative-heading-rank; base=h2

\section{第八首}

（应答—— \Emph{你们有福了，你们的身体已被照亮！}）\par

1. 天主以祂的仁慈俯身降下——将祂的怜悯与水调和——并将祂威严的本性——与人类可怜的身体结合。——祂藉水制造时机——降下并居住在我们内：——如同仁慈的时机——祂降下并居住在胎中：——啊，天主的仁慈——祂为自己寻找一切时机居住在我们内！\par

2. 祂俯身降下到\Place{曷勒布}的洞穴——使祂的威严居住在\Person{梅瑟}身上——祂将荣耀的光辉分施给必朽之人。——那里有圣洗的预象：——祂降下并居住在其中的——在水中调和——祂威严的大能——以便居住在软弱者内。——气息停留在\Person{梅瑟}身上，——而你们身上却有\OriginalTerm{基督的成全}{Perfecting of Christ}。\par

3. 那大能当时无人能承受——连\Person{梅瑟}，拯救者之首也不能，——连\Person{厄里亚}，热诚者之首也不能；——\Person{色辣芬}也遮掩自己的面容——因为那是制服万有的大能。——祂的仁慈将柔和——调和在水和油中——好使软弱的人类——能在祂面前站立——被水和油所遮盖。\par

4. 被掳的司铎们又在井中——藏匿并隐藏了圣所的火，\BibleRef{加上 1:19}——那是那荣耀之火的奥秘——\Person{大司祭}将其调和在圣洗中。——司铎们取了淤泥——撒在祭台上——因为那井的火——已与淤泥混合——这是我们身体的奥秘，在水中——与圣神的火混合了。\par

5. 在\Place{巴比伦}的那三位著名人物——在火窑中受了洗，然后出来——他们进去，在火焰的洪流中沐浴，被炽热的波涛拍打。——那里有从天降下的露水——洒在他们身上——解开了他们身上——属地君王的锁链。——看！那三位著名人物进去，在火窑中找到了第四位。\par

6. 那外在得胜的可见之火——指向圣神的火——它已调和，看！隐藏在水中。——在火焰中预象了圣洗，——在那窑的烈焰中。——来吧，进去，受洗吧，我的弟兄们——因为看！它解开了锁链；——因为其中居住并隐藏着——\OriginalTerm{天主的调解人}{Daysman of God}——祂在火窑中是第四位。\par

7. 我们的主又说了两句话——它们以同一声音和谐一致：——祂说：\Quote{我来是要送火}，——又说：\Quote{我应当受一种洗。}——藉着圣洗的火，恶者所点燃的火——被熄灭了：——而圣洗的水克服了——那些争闹的水——恶者曾用它们试探——\Person{若瑟}，他得胜并戴上了冠冕。\par

8. 看！我们救赎主的纯火——祂以仁慈在人类中点燃！——藉着祂的火，祂熄灭了那火——那在污秽和罪人中所点燃的火。——这是那火，荆棘——和莠子在其中被烧尽。——但你们的身体有福了——它们已在火中受洗——这火烧尽了你们的荆棘——而你们的种子藉此发芽升向天堂！\par

9. \Person{耶肋米亚}在母胎中被圣化并受教。——但如果婚姻的低微怀抱——在怀孕和生产他时被圣化——那么圣洗岂不更要圣化——那些纯洁而属神之人的——怀孕和生产！——因为在那里，在母胎内——是所有人的怀孕——但在这里，从水中——是属神之人所配得的出生。\par

10. 因为\Person{耶肋米亚}虽然在母胎中被圣化——他们却拿起钉子，将他扔进坑中。——先知在他被玷污中仍是圣的，——因为他的心洁净，虽在淤泥中。——要惧怕，我的弟兄们——因为看！今天你们隐秘的玷污——和你们罪恶的憎恶——被洗去了。——不要再转向不洁——因为你们的身体只有一次洁净！\par

11. 那受洗后又犯罪的狂妄者——如同蛇脱了皮又穿上，更新而变得年轻，却又转回——重新穿上它旧日的皮——因为蛇并没有——脱去它的本性。——要脱去那诱惑者——灵魂的败坏者——就是\OriginalTerm{旧人}{old man}——不要让那新穿上的——新衣变旧！\par

12. \Person{厄里叟}将木头丢进水中，使重的漂浮、轻的下沉：——它们的本性在水中交换了。——那里发生了一件新事，不按本性。——那么，主啊，这对祢的恩宠更容易：使沉重之罪——在水中下沉——而轻的灵魂——被拉上来升到高处！\par

13. \Person{若苏厄}对\Place{耶里哥}——咒诅了它的城墙，也咒诅了它的泉源。——被\Person{若苏厄}咒诅以至于灭亡的——在耶稣的奥秘中又蒙祝福。——有人将盐撒在其中——它们就愈合了、变甜了：——这盐的奥秘——那来自\Person{玛利亚}的甜盐——调和在水中——藉此我们瘟疫的恶臭得了医治。\par

14. 看！平静的水在你面前——圣洁、安宁而怡人——因为它们不是争闹的水——将\Person{若瑟}投进地牢——也不是那些争闹的水——百姓在旷野旁——争闹并抱怨。——有一些水——使人与天和好。\par

15. \Person{哈加尔}看见了一眼泉水——从中给她那野性的儿子饮水，他在旷野中成了野驴。——那泉水的代替者是圣洗。——\Person{哈加尔}的众子在其中受洗——变得温和而平安。谁曾见过这样的公羊——它们被套上轭，看！劳苦——与驯良的小公牛一起——它们的耕种的种子收获百倍！\par

16. 起初，那运行在水面上的\Person{圣神}——\newline{}使它们受孕并生出——\newline{}蛇、鱼和鸟。——\newline{}\Person{圣神}已在\OriginalTerm{圣洗}{Baptism}中运行，——\newline{}并在奥秘中生出了鹰——\newline{}贞女与主教；——\newline{}并在奥秘中生出了鱼——\newline{}独身者与代祷者；并在蛇的奥秘中，——\newline{}看！狡猾的已变得像鸽子般纯朴！\par

17. 看！我们主的剑在水中！——\newline{}那把儿子与父亲分开的：——\newline{}因为它是活剑，使——\newline{}活人与死人分离，看！——\newline{}看！他们受了洗，成为——\newline{}贞女与圣人——\newline{}那些下去、受了洗、穿上——\newline{}独一圣子的人。——\newline{}看！许多人已大胆地来到他面前！\par

18. 因为凡受了洗并穿上他——\newline{}独生者、众多之主——\newline{}就因此充满了众多之位——\newline{}因为基督成了他的巨大宝藏：——\newline{}因为他在旷野中成了——\newline{}美味之席——\newline{}他在婚宴上成了——\newline{}精选美酒之泉。——\newline{}他对众人成为\Emph{如此}在一切事上——\newline{}藉助、医治和应许。\par

19. 厄里叟与守望者同等——\newline{}在他的作为中，荣耀而圣洁。——\newline{}守望者的营帐围绕他——\newline{}因此让\OriginalTerm{圣洗}{Baptism}为你成为——\newline{}守护者的营帐，——\newline{}因为藉着它，心中住着——\newline{}下面之人的希望——\newline{}和上面之人的主。——\newline{}为他的缘故圣化你们的身体——\newline{}因为他所住之处，腐败不能靠近。\par

20. 那海的水已不复存在——\newline{}它曾以巨浪保护了百姓——\newline{}又以巨浪击倒了众民。——\newline{}\OriginalTerm{圣洗}{Baptism}中的水有相反的功效。——\newline{}在其中，看！百姓得生命——\newline{}在其中，看！百姓灭亡：——\newline{}因为凡未受洗者——\newline{}在赋予众人生命的水中——\newline{}他们隐形地死去。\par

21. 那海的水已不复存在——\newline{}它们曾暴怒，向约纳翻腾，——\newline{}将阿米泰的儿子投入深渊。——\newline{}他虽逃跑，却被囚在监牢中——\newline{}\Emph{天主}将他投入并捆绑——\newline{}在深牢中的深牢——\newline{}因为他在海中捆绑了他——\newline{}又在鱼中捆绑了他。——\newline{}恩宠为他作保——\newline{}她开了监牢，带出了宣讲者。\par

22. 先知们称至高者为火——\newline{}\Quote{吞噬之火，}和\Quote{谁能与之同住？} \EditorialNote{见：\BibleRef{依 30:27}}——\newline{}百姓不能住在其中——\newline{}它的力量压垮了众民，他们惊惶。——\newline{}在其中，你已受傅油；——\newline{}你在水中穿上了他——\newline{}在面包中你吃了他——\newline{}在酒中你喝了他——\newline{}在声音中你听到了他——\newline{}并在心眼看见了他！\par

% structure: p0214 source=h2 target=section reason=relative-heading-rank; base=h2

\section{赞美诗 9}

(应答，\Emph{降来并圣化水以赦免亚当子孙之罪的，当受赞美！})\par

1. 若翰，你看见了圣神——\newline{}停留在圣子的头上——\newline{}以显示至高的头——\newline{}下去受了洗——\newline{}上来成为地上的头！——\newline{}你们因此成了圣神的子女——\newline{}基督成了你们的头：——\newline{}你们也成了他的肢体。\par

2. 请思想并看你们何其崇高——\newline{}你们不再有约旦河——\newline{}而是荣耀的\OriginalTerm{圣洗}{Baptism}，其中有平安——\newline{}展开双翼荫庇你们的身体。——\newline{}若翰在旷野中施洗：——\newline{}在圣洗的纯净洪流中，——\newline{}你们在其中被纯净地施洗。\par

3. 婴儿们看见它的荣耀时以为——\newline{}它的权能因它的辉煌而增强。——\newline{}但它本身不变，在自身之内——\newline{}并不分裂。——\newline{}但从来不减少或增加的权能——\newline{}在我们中是小的，或是大的，——\newline{}而多有智慧的人——\newline{}\OriginalTerm{圣洗}{Baptism}在他中也是大的。\par

4. 人的知识，若被高举——\newline{}他的位置也高于他的弟兄；——\newline{}而他的信德若大——\newline{}\Emph{同样}他的应许也大——\newline{}而他的智慧怎样，他的冠冕也怎样。——\newline{}光虽全然美好——\newline{}且自身完全相等——\newline{}\Emph{然而}一只眼比另一只更美好。\par

5. 耶稣将他的权能混合在水中：——\newline{}我的弟兄们，作为有辨识力的\Emph{人}，穿上他！——\newline{}因为有些人在水中——\newline{}只感知自己被洗了。愿我们的灵魂与身体同被洗！——\newline{}让身体感知可见的水——\newline{}灵魂感知隐秘的权能——\newline{}好使你们在可见与隐秘的事上相似他！\par

6. \OriginalTerm{圣洗}{Baptism}何其美丽——\newline{}在心的眼中；来，让我们凝视它！——\newline{}如同被印玺塑造；——\newline{}接受它的形象——\newline{}好使我们的形象无所缺乏！——\newline{}因为心洁白的羊群——\newline{}凝视水中的荣耀：——\newline{}在你们的灵魂中反射它！\par

7. 水本性地如明镜——\newline{}使人在其中审视\Emph{自己}。——\newline{}振作你的灵魂，有辨识力的你——\newline{}并与之相似！——\newline{}因为它在其中映出你的形象——\newline{}从它、在它身上找到榜样——\newline{}在其中凝视\OriginalTerm{圣洗}{Baptism}——\newline{}并穿上其中隐藏的美！\par

8. 对那听见声音——\newline{}却不知其意义的人有何益处？——\newline{}谁听见声音却没有——\newline{}对它的理解——\newline{}他的耳朵充满但灵魂空虚。——\newline{}看！既然恩赐丰盛——\newline{}就请以辨识力接受它。\par

9. 带有理解的\OriginalTerm{圣洗}{Baptism}——\newline{}是两种光的结合——\newline{}其光线的源泉丰富。——\newline{}......——\newline{}心智上的黑暗离去——\newline{}灵魂在美中看见他——\newline{}隐藏的荣耀基督——\newline{}当\Emph{荣耀}逝去时便忧伤。\par

10. 没有理解的\OriginalTerm{圣洗}{Baptism}——\newline{}是充满却空虚的宝藏——\newline{}因为接受它者在其中贫乏——\newline{}因为他不知道——\newline{}他进入并居住的财富何其大。——\newline{}因为其中的恩赐大——\newline{}虽然卑微的人不觉察——\newline{}自己被高举如同它一样。\par

11. 开阔你们的心智，我的弟兄们——\newline{}看空中的秘密柱子，其根基从水中央——\newline{}直抵至高之处的门，如同雅各伯看见的梯子。——\newline{}看！藉此光降临\OriginalTerm{圣洗}{Baptism}，——\newline{}\Emph{藉此}灵魂升上天堂——\newline{}好使我们在一爱中相融。\par

12. 我们的主受若翰洗礼时——\newline{}发出十二道泉源——\newline{}它们涌流，用其溪流洗净——\newline{}众民的污秽。——\newline{}他的朝拜者被洗白如他的衣服——\newline{}大博尔山上的衣服和水中的身体。——\newline{}众民代替衣服被洗白——\newline{}并成了他荣耀的衣袍。\par

13. 从你们的衣服上学习，我的弟兄们，——你们的肢体应当如何保持。——因为衣服，虽然可以多次洗净——仍为它的美观而妥善保存——那仅有一次圣洗的身体——它的保存更应倍加谨慎——因为它的危险甚多。\par

14. 太阳进入狭窄的房屋——虽大却被局限其中；——但在美好宽广的房屋——当它升起其上——四面八方放射光芒——太阳虽同一本性——在不同房屋中却经历变化：——我们的主在不同的人中也是如此。\par

% structure: p0230 source=h2 target=section reason=relative-heading-rank; base=h2

\section{圣诗 10}

（对经。——\Quote{光荣归于来临并恢复它的那一位！}）\par

1. 亚当犯罪，招致一切忧苦——世界照样，染上一切罪债。——它不思如何得以恢复——反思堕落如何令它愉悦。——光荣归于来临并恢复它的那一位！\par

2. 这原因召叫那位纯洁者——祂应来受洗，与污秽者一同。——天为祂的荣耀而裂开。——为使万有的净化者与众人一同受洗——祂降来，圣化我们圣洗的水。\par

3. 祂为何进入母胎——也为同一原因进入河流。——祂为何进入坟墓——也为同一原因使我们进入祂的内室。——祂为每一原因成全了人类。\par

4. 祂的成胎是我们福泽的宝库；——祂的诞生是我们喜乐的宝藏；——祂的圣洗是我们赦罪的缘由；——祂的死亡是我们生命的缘由。——唯独祂在复活中战胜了死亡。\par

5. 祂诞生时，光明之星闪耀于空中；——祂受洗时，光明从水中射出；——祂死亡时，太阳在穹苍昏暗；——祂受难时，众星与祂一同隐没；——祂显现时，众星与祂一同升起。\par

6. 祂的荣耀因祂的尊威而显明；——祂的苦难因祂的人性而显明；——祂的慈爱因祂的恩惠而显明；——祂的审判因祂的公义而显明。——祂将祂的属性倾注于属祂的人。\par

7. 使那曾见祂荣耀而藐视祂的人——再次瞻仰祂的荣耀而朝拜祂；——使那曾不屑品尝祂恩惠的人——恐惧自己会感受到祂的公义。——祂将祂的帮助倾注于敬拜祂的人。\par

8. 看！东方在清晨被照亮！——看！南方在正午变黑暗！——西方又在傍晚被照亮。——三个方向代表同一个诞生；——它们宣告祂的死亡和生命。\par

9. 祂的诞生涌流，与祂的圣洗相连；——祂的圣洗又涌流，直至祂的死亡；——祂的死亡引领并达到祂的复活——一座四重桥梁通往祂的国度；看！祂的羊群跟随祂的足迹走过。\par

10. 正如除非经过出生的门——无人能进入受造界；——同样，除非经过复活的門——无人能进入天国；——而那切断自己桥梁的人，便毁灭了自己的希望。\par

11. 祂穿上铠甲，得胜并加冕；——祂将铠甲留在地上，升了天；——若有人渴望冠冕，——他当取这铠甲，借此赢得——他所渴求的胜利冠冕。\par

12. 祂在地上成就了义，并升了天。——但若那洗净万有的祂受了洗——还有谁人不该受洗呢？——因为恩宠已临到圣洗——洗去我们伤口的污秽。\par

13. 天主的强迫是无可抗拒的力量——[但那出于强迫的并不中悦祂，]——正如出于明辨意志的。——因此祂藉我们的果实召叫我们——我们不是如同受强迫而活，却是出于劝导。\par

14. 祂是良善的，看！祂在这两件事上劳苦——祂不愿强迫我们的自由——也不容许我们滥用它。——因为若祂强迫，就剥夺了它的能力——若祂放任，就使它失去帮助。\par

15. 祂知道若祂强迫，便剥夺我们；——祂知道若祂抛弃，便毁灭我们；——祂知道若祂教导，便赢得我们。——祂没有强迫，也没有抛弃，如那恶者所为：——祂教导、惩戒并赢得了我们，如良善的天主。\par

16. 祂知道祂的宝库丰盈：——祂将宝库的钥匙交在我们手里。——祂使十字架成为我们的司库——为我们开启乐园之门，——正如亚当开启了地狱之门。\par

% structure: p0248 source=h2 target=section reason=relative-heading-rank; base=h2

\section{圣诗 11}

（对经。——\Quote{让那被恶者剥去衣裳的身体欢喜，因它们在水里穿上了荣耀！}）\par

1. 女儿啊，感恩吧，因你的加冕加倍——看！你的圣殿和你的儿子一同喜乐。——你圣殿的奉献在于礼仪——你儿子的奉献在于傅油。——你是有福的，同时……——……为住在你内的人搭设帐幕——圣神也住在你儿子身上！\par

2. 我们的主开启了圣洗——在约旦河这条蒙福的河中。——高处和深处都因祂而喜乐。——祂从水中带出祂和平的初果——因为它们是初果，圣洗的果实。——良善的天主以祂的慈悲成就——使和平成为地上的初果。\par

3. 梅瑟支搭了暂时的会幕；——司祭们用水沐浴——进去供职；却又被打伤受罚——因他们内心没有洁净。——你有福了，因在伟大苦难的逾越节里——司祭们藉他们祭献的馨香——看！正在你内洁净灵魂！\par

4. 先知所见的是伟大的奥秘——那汹涌的洪流。——他凝视其深处，看到了——你的美丽，而非他自己；他见到的是你，因你的信德永存——你那无形的洪流将淹没——偶像崇拜的诡计。\par

5. 虽然若翰在妇人所生中最大——但较小的他却比他更大——因为他的受洗者再次受洗——在宗徒的圣洗中。——你有福了，因你的司祭比他更大——仅此一点，因他的圣洗——永远存留。\par

6. 那来自\Place{史罗亚}的圣洗——并未给躺在\Emph{那里}的人带来仁慈——他等候了三十八年——因为他是肋未人的徇私者。——\Emph{你可赞颂}，因你的医治\Emph{在}你内为众人具备——你的司铎们虔诚而准备——为一切\Emph{需要}你帮助的人。\par

7. 先知治愈了不洁的水——医治了荒芜之地的病患——以致其死亡被除去，其区域回荡，因其后代增多，其腹中充满。——\Emph{主}，您的恩宠大于\Person{厄里叟}的！——请在我泉源的大河中增多我的羔羊和羊群！\BibleRef{则 47:1}\par

8. 在你居所内的奇迹何其伟大——羊群与牧者一起——在水流旁的那些——两位不可见与一位可见的施洗者。——在它们的泉源中受圣洗者\Emph{是}有福的！——因三只手臂扶持了他——三个名字保护了他！\par

% structure: p0258 source=h2 target=section reason=relative-heading-rank; base=h2

\section{第十二首讚歌}

（应答.— \Emph{祂下降并在约但河领受圣洗，使百姓从谬误中回转，祂是可赞美的！}）\par

1. 在圣洗中，\Person{亚当}\Emph{再次}找到了——那在\Emph{伊甸园}树木间的荣耀。——他下降，从水中领受了它——他穿上它，上升并因此装饰。——愿怜悯众人的祂受赞颂！\par

2. 人在乐园中堕落了——圣洗中的慈悲恢复了他：——他因\Emph{撒殚}的嫉妒失去了其俊美——又因\Emph{天主的}恩宠\Emph{再次}找到了它。——愿怜悯众人的祂受赞颂！\par

3. 那对夫妇在伊甸园中被装饰——但蛇偷走了他们的冠冕：——然而慈悲碾碎了那受诅咒者——并使那对夫妇在他们的衣裳中光鲜。——愿怜悯众人的祂受赞颂！\par

4. 他们以需要的叶子蔽体；——但那慈悲者怜悯了他们的美丽——并以树木的叶子代替——祂在水中以荣耀给他们穿上。——愿怜悯众人的祂受赞颂！\par

5. 圣洗是生命的泉源——\Person{天主之子}以祂的生命开启——并从祂的肋旁涌出了溪流。——凡口渴的，来！来，欢欣！——愿怜悯众人的祂受赞颂！\par

6. 圣父已封印\Emph{圣洗}，为高举它——圣子已迎娶它，为光荣它——圣神以三重封印——盖印其上，它便闪耀在圣洁中。——愿怜悯众人的祂受赞颂！\par

7. 不可探究的圣三——在圣洗中贮藏了宝藏。——贫穷的人，下到它的泉源！——困乏的人，从中致富！——愿怜悯众人的祂受赞颂！\par

% structure: p0267 source=h2 target=section reason=relative-heading-rank; base=h2

\section{第十三首讚歌}

\Emph{受圣洗者的讚歌。}\par

（应答.— \Emph{弟兄们，请讴歌赞颂，向万物之主的圣子；祂为你们繫上了王侯所渴慕的冠冕！}）\par

1. 你们的衣裳如雪闪亮，我的弟兄们——你们的光辉美丽无比，如同天使一般！\par

2. 如同天使，你们上来了，亲爱的——从约但河的河中，带着圣神的武装。\par

3. 那不败坏的婚房，我的弟兄们，你们领受了：——今天你们穿上了\Person{亚当}之家的荣耀。\par

4. 那\Emph{出于}果子的审判，是\Person{亚当}的定罪：——但今天为你们升起了胜利。\par

5. 你们的衣着闪亮，你们的冠冕美好：——今天首生者藉司铎之手为你们繫上了。\par

6. \Person{亚当}在乐园中领受了祸患：——但今天你们领受了荣耀。\par

7. 胜利的武装，你们穿上了，我亲爱的：——在司铎呼求圣神的时候。\par

8. 天使欢欣，地上的人雀跃：——在你们的庆节中，我的弟兄们，其中毫无污秽。\par

9. 天堂的美物，我的弟兄们，你们已领受：——当心那恶者，免得他掠夺你们。\par

10. 天上君王出现的那一天：——祂为你们敞开祂的门，并邀请你们进入伊甸园。\par

11. 不衰败的冠冕，戴在你们头上：——让你们的口时时歌颂赞美诗。\par

12. \Person{亚当}因着果子，\Emph{天主}忧伤地驱逐了他：——但祂在喜乐的婚房中使你们欢喜。\par

13. 谁不会在你们的婚房中欢欣呢，我的弟兄们？——因为圣父、圣子与圣神在你们内欢欣。\par

14. 圣父要做你们力量的墙壁：——圣子做救赎者，圣神做护卫。\par

15. 殉道者藉自己的血光荣了他们的冠冕：——但我们的救赎者藉祂的血光荣了你们。\par

16. 守夜者和天使，因悔改者而喜乐：——他们必因你们而喜乐，我的弟兄们，因为你们与他们相似。\par

17. \Person{亚当}在乐园中没有尝过的果子：——今天已喜乐地放在你们口中。\par

18. 我们的救赎者以树预表了祂的身体：——\Person{亚当}没有尝过，因为他犯了罪。\par

19. 恶者争战并征服了\Person{亚当}之家：——藉着你们的圣洗，我的弟兄们，看！今天他被征服了。\par

20. 胜利是伟大的，但今天你们赢得了：——如果你们不疏忽，你们就不会灭亡，我的弟兄们。\par

21. 光荣归于被披上的人，光荣归于\Person{亚当}之家！——在那\Emph{是}来自水中的新生中，让他们欢欣并蒙福！\par

22. 赞美那使祂的教会穿上荣耀的祂！——光荣那高举\Person{亚当}之家族的祂。\par

% structure: p0292 source=h2 target=section reason=relative-heading-rank; base=h2

\section{第十四首讚歌}

\Emph{论及我们主与\Person{若翰}的讚歌。}\par

（应答.— \Emph{光荣归于祢}，我主，因为祢\Emph{——天地喜乐地朝拜！}）\par

1. 我的思绪带我到了\Place{约但河}——我看见一个奇观，当时那荣耀的新郎显现——祂要给新娘带来自由和圣洁。\par

2. 我看见\Person{若翰}充满惊奇——群众站在他周围——那荣耀的新郎俯首——向那不育者的儿子，好让他给祂施洗。\par

3. 面对圣言和声音，我的思绪惊奇：——因为看！\Person{若翰}就是那声音；——我们主显明为圣言，使那隐藏的得以彰显。\par

4. 新娘已订婚却不知道——她凝视的新郎是谁：——宾客聚集，旷野充满——我们主隐藏在他们中。\par

5. 那时新郎显示了自己——祂向\Person{若翰}的声音靠近：——那前驱感动了并论祂说——\Quote{这就是我所宣告的新郎。}\par

6. 祂来领受圣洗，祂是给众人付洗的那位——并在\Place{约但河}显示了自己。——\Person{若翰}看见祂便后退——恳求地，于是他这样说：——\par

7. \Quote{我主，祢如何愿意领受圣洗——祢在祢的圣洗中赎免一切？——圣洗仰望祢——祢可向它倾注圣洁和完美吗？}\par

8. 我们的主说：\Quote{我愿\Emph{如此}——上前来，给我付洗，好使我的旨意得以承行。——你们不能反抗我的旨意：——我要由你受洗，因为我愿\Emph{意}这样。}\par

9. \Quote{我恳求，我主，不要强迫我——因祢向我说的话实在严厉——\InnerQuote{我需要你给我付洗；}——因为祢是用祢的牛膝草洁净万有的那一位。}\par

10. \Quote{我已要求，且乐意如此成就——而你，若翰，为何反对？——让义德得以圆满——来，给我付洗；你为何站在那里？}\par

11. \Quote{人怎能公然用手——握住燃烧的火？——\Emph{祢是}火，求祢怜悯我——不要命我靠近祢，因我实在为难！}\par

12. \Quote{我已向你显明我的旨意；你还追问什么？——上前来，给我付洗，你必不被焚烧。——洞房已备妥，不要阻止我——赴那已准备好的婚宴。}\par

13. \Quote{守望者恐惧，不敢——注视祢，以免失明——而我，我主，怎能给祢付洗？——我软弱无力上前，请不要责怪我！}\par

14. \Quote{你既害怕，就不要反对——我所渴望的旨意：——洗礼原是为了我。——完成你所蒙召的工程吧！}\par

15. \Quote{看！我在约但河曾向你宣告——在那些不信的百姓耳中——若他们看见祢由我受洗——必怀疑祢不是主。}\par

16. \Quote{看！我\Emph{要}在他们眼前受洗——那派遣我的父必为我作证——我是祂的子，祂所喜悦的——为使那在\Emph{祂}义怒之下的亚当得以和好。}\par

17. \Quote{我主，我理当认识我的本性——我是由泥土所造——而祢是塑造万有的陶工——那么，我为何要给祢用水付洗？}\par

18. \Quote{你理应知道我为何而来——以及我为何要你为我付洗。——这是我行走之路的中途——不要阻止洗礼。}\par

19. \Quote{祢所到的河流何其窄小——祢竟愿栖息其中，让它洁净祢。——诸天不足容纳祢的伟大——何况洗礼怎能容纳祢！}\par

20. \Quote{母胎比约但河更小——然而我曾甘愿居于童贞女之中——正如我从妇人所生——同样\Emph{也要}在约但河受洗。}\par

21. \Quote{看！万军肃立！——守望者的队列，看！他们朝拜！——若我上前，我主，给祢付洗——我战栗恐惧不已。}\par

22. 万军与群众都称你有福——全都因你为我付洗。——为此我从母胎中拣选了你——不要怕，因我已愿\Emph{意}如此。\par

23. \Quote{我既受派遣，已预备道路——我既受命，已聘定新娘。——愿祢的显现传遍普世——如今祢既来了，就不要让我给祢付洗！}\par

24. \Quote{这是我的准备，因我如此愿意——我要下去在约但河受洗——为那些受洗者磨亮盔甲——使他们在我内洁白，而我不被征服。}\par

25. \Quote{圣父之子，我为何要给祢付洗？——看！祢在父内，父在祢内。——祢赐予司铎圣洁——为何求取这普通的水？}\par

26. \Quote{亚当的子孙仰望我——要我为他们成就新生。——我要在水中为他们寻出一条路——若我不受洗，\Emph{这事}便不能成就。}\par

27. \Quote{祢的\OriginalTerm{大司祭}{Pontiffs}被祝圣——司铎由祢的牛膝草得以洁净——祢也立了受傅者和君王。——洗礼如何能有益于\Emph{祢}？}\par

28. \Quote{你所聘定的新娘正等候我——要我下去受洗，祝圣她。——新郎的朋友，不要阻止我——前往那等候我的\OriginalTerm{洗礼}{washing}。}\par

29. \Quote{我不能，因我软弱——无法将祢的火焰握在手中。——看！祢的军团如火焰——派一个守望者来给祢付洗吧！}\par

30. \Quote{我的身体并非取自守望者——以致我要召一个守望者来为我付洗。——看！我已穿上亚当的身体——而你，亚当之子，理当为我付洗。}\par

31. \Quote{诸水看见祢，就大大惧怕；——诸水看见祢，看！它们颤抖！——河因恐惧而翻腾——我\Emph{是}软弱的，怎能给祢付洗？}\par

32. \Quote{诸水在我洗礼中被祝圣——它们将从我领受火与圣神——若我不受洗，它们就不能得以成全——不能结出不朽的子女。}\par

33. \Quote{火若靠近祢的火——必如碎秸被烧尽。——西乃山尚且不能承受祢——我\Emph{是}软弱的，要在哪里给祢付洗？}\par

34. \Quote{我是烈火——但为人的缘故，我成为婴孩——在童贞女的胎中。——如今我将在约但河受洗。}\par

35. \Quote{祢为我付洗甚为合宜——因祢有圣洁能洁净一切。——污秽者在祢内得以成圣；但祢\Emph{是}圣洁的，为何要受洗？}\par

36. \Quote{你为我付洗甚为合宜——照我所命，不要反对。——看！我在母胎中已为你付洗——如今你在约但河为我付洗吧！}\par

37. \Quote{我是奴隶，又是软弱的。——释放众人的祢，求祢怜悯我！——我不配解祢的鞋带——谁能使我有幸\Emph{触摸}祢崇高的头？}\par

38. \Quote{奴隶们在我的洗礼中获得释放——罪状在我的洗涤中被涂抹——释放令在水印中被封印；——若我不受洗，这一切便归乌有。}\par

39. \Quote{空气披着一件火披风——在约但河上空等候祢；——若祢同意并愿意受洗——祢就自己付洗，完成一切。}\par

40. \Quote{理应由你为我付洗——免得有人误解我，说：——\InnerQuote{祂若非与父家疏远——为何肋未人惧怕为祂付洗？}}\par

41. \Quote{那么，当祢受洗时——我该在约但河上如何完成祷词？——当圣父与圣神显现于祢之上时——作为司铎，我该呼求谁？}\par

42. \Quote{祷词当在静默中完成：——来，只将你的手按在我身上。——圣父将代替司铎——说出关于祂圣子合宜的话。}\par

43. \Quote{被请的人，看，他们都站着——新郎的宾客，看，他们作证——我每日在他们中间说——\InnerQuote{我是声音，不是圣言。}}\par

44. \Quote{在旷野呼喊者的声音——完成你来应做的工作——好使你前往的旷野——回荡你所宣讲的浩大平安。}\par

45. \Quote{守望者的呼声已传入我耳——看，我从圣父的家听见——发出呐喊的万军——\InnerQuote{新郎，在你的主显节，诸世界获得了生命。}}\par

46. \Quote{时间紧迫，婚宴的宾客——看着我，看正在发生的事。——来，给我施洗，好使他们赞美——当圣父的声音被听见时！}\par

47. \Quote{我听从，我的主，照你的话：——因你的爱催迫你，前来受洗！——尘土敬拜他所达到的——就是他应该亲手按在塑造他的那一位上。}\par

48. 天上的众军肃立静默——新郎下到约旦河；——圣者受了洗，立刻上来——祂的光辉照耀世界。\par

49. 至高之处的门被打开——听见了圣父的声音——\BibleQuote{这是我的爱子，我所喜悦的。}——万民啊，来朝拜祂。\par

50. 看见的人都惊奇地站着，因降临并为他作证的圣神。——赞美你的主显节，它使万有喜乐——你的启示光照了诸世界！\par

% structure: p0345 source=h2 target=section reason=relative-heading-rank; base=h2

\section{赞美诗 15}

1. 在圣子诞生时，光明破晓，——黑暗从世界逃遁，——大地被光照；愿它因此将荣耀——归于照亮它的圣父的光明！\par

2. 祂从童贞玛利亚的胎中升起，——当祂显现时阴影消逝，——错误的黑暗被祂扼杀，——地极被光照，好将荣耀归给祂。\par

3. 万民之中有大骚动——在黑暗中光明破晓，——万邦欢欣，将荣耀归于祂，因在祂的诞生中他们都得了光照。\par

4. 祂的光照耀东方——波斯被星光所照：——\Emph{祂的}主显节向她传报佳音并邀请她——\Quote{祂是为带来喜乐的祭献而来。}\par

5. 光明之星急忙前来，穿越黑暗破晓，并召唤他们——好让万民前来，在降临大地的伟大光明中欢腾。\par

6. 穹苍从众星中派遣一位使者——向他们宣告——向波斯子孙宣告，好让他们预备——迎接君王并朝拜祂。\par

7. 大亚述察觉\Emph{此事}——便召唤贤士并对他们说——\Quote{带上礼物去吧，尊敬祂——在犹大升起的大君王。}\par

8. 波斯的首领们欢欣鼓舞——携带来自本地的礼物——他们带给童贞玛利亚之子——黄金、没药和乳香。\par

9. 他们进入，见祂是个婴孩——住在卑微妇女的家中——他们上前欢然朝拜——将宝物呈在祂面前。\par

10. 玛利亚说：\Quote{这些是为谁？——有什么用途？是什么原因——促使你们带着宝物从本国——来到这孩子面前？}\par

11. 他们说：\Quote{你的儿子是君王——祂戴冠冕，是万有的君王——祂对世界的权能是\Emph{伟大}的——万有都要服从祂的王国。}\par

12. \Quote{这事何时发生——卑微的妇女竟产下君王？我本是贫乏穷困——君王怎会从我而出？}\par

13. \Quote{唯独在你身上这事成就——一位大能的君王要自你出现——在你身上贫穷要显扬——众王冠都将向你的儿子臣服。}\par

14. \Quote{我没有君王的财宝——财富从未降临于我。——我的房屋卑微，住处穷乏——为何你们宣告我儿子是君王？}\par

15. \Quote{你的儿子有大宝藏——有足以使众人富足的财富；——因为君王的财宝会耗尽——但祂永不匮乏，不可测量。}\par

16. \Quote{或许另有一位为你们而生——那位君王，你们去查问吧。——这只是卑微妇女的儿子——一个\Emph{不配觐见君王的人。}}\par

17. \Quote{光明岂会迷失——所差遣的道路？不是黑暗召唤并引领我们——我们在光中行走，你的儿子是君王。}\par

18. \Quote{看，你见一个不会说话的婴孩——祂母亲的家空空如也，贫乏不堪——里面没有君王之物——那么，怎能从中看见君王呢？}\par

19. \Quote{看！我们看见那君王无言安息，\Emph{是}卑微如你所说：——但再次我们看见众星——在高天，他命令它们急忙传报祂。}\par

20. \Quote{人啊，你们应当先查问——谁是君王，然后才朝拜他——免得\Emph{你们的}道路或许错了——而另有其人是那诞生的君王。}\par

21. \Quote{少女啊，你应当接受——我们已学到你的儿子是君王——从这光明之星，它从不差错——道路平坦，他已引领了我们。}\par

22. \Quote{这孩子还小，看！他没有——君王的冠冕和宝座——你们看见了什么，竟以宝物向他致敬——如同对待君王？}\par

23. \Quote{他是幼小的，因为祂为安详的缘故而愿意\Emph{如此}——如今谦和，直到祂被启示。——时候将到，那时一切冠冕——都要俯伏朝拜祂。}\par

24. \Quote{他没有军队——我儿子也没有军团和队伍：——祂住在母亲的贫穷中——为何你们称祂为王？}\par

25. \Quote{你儿子的军队在高天——他们乘车在高处，燃烧如火——其中一位前来召唤我们——我们全国都惊慌。}\par

26. \Quote{这孩子是婴孩，怎可能——祂是君王，而不为世人所知？——那些有能有名的人——一个婴孩怎能统治他们？}\par

27. \Quote{你的婴孩是年长的，童贞玛利亚——万古常存者，超越万有，亚当在他旁边简直是婴孩——在祂内，\Emph{万有}都成为新的。}\par

\Quote{你应当陈明——全部奥秘，并解释\Emph{它}——那位\Emph{它}向你启示我圣子奥秘者是谁——祂是你地域的君王。}\par

\Quote{同样，你应当接受这事——若非真理引领了我们，我们不会从地极流浪至此——也不会为你圣子而来。}\par

\Quote{那在你们那里所行的一切奥秘——现在就以朋友的身份向我揭示——是谁召唤你们到我这里来？}\par

\Quote{一颗大星向我们显现——极其荣耀，远超众星——我们的地因它的火光被点燃——它向我们传报，这位君王已出现。}\par

\Quote{我恳求你们，不要在我们境内谈论这些——免得他们发怒，\Emph{并}地上的众王联合起来——因嫉妒敌视这婴孩。}\par

\Quote{不要惊慌，童贞女！——你的圣子将粉碎一切冠冕，将它们置于祂的脚下——他们不能征服他们所嫉妒的那一位。}\par

\Quote{我因黑落德而害怕——那污秽的狼，恐怕他袭击我——拔出剑来，用它砍断——那未成熟的甜美葡萄串。}\par

\Quote{不要因黑落德而害怕——因为他的宝座已安置在你圣子手中——一旦祂开始统治，宝座将被夷平——他的冠冕将坠于地上。}\par

\Quote{耶路撒冷是血河——其中卓越者被杀害——若她察觉祂，必会攻击祂。——要隐秘地讲，不要将\Emph{它}张扬出去。}\par

\Quote{一切激流，同样刀剑——都必因你圣子之手平息——耶路撒冷的剑将变钝——完全不再渴望杀戮。}\par

\Quote{耶路撒冷司铎的经师们——倾流鲜血，毫不在意。——他们将激起杀伐的争端——反对我和这婴孩；贤士们，请静默！}\par

\Quote{经师和司铎们将无法——因嫉妒伤害你的儿子——因为藉着祂，他们的司铎职将被解散——他们的节期化为虚无。}\par

\Quote{一位\OriginalTerm{守望者}{Watcher}向我启示，当我领受——这婴孩的成孕时，我儿子是君王——祂的冠冕来自高处，永不消解——祂向我宣告，正如你们所\Emph{做}的那样。}\par

\Quote{因此，你所说的那位守望者——就是化作星而来的那位——向我们显现，并为我们带来喜讯——祂伟大而荣耀，超越众星。}\par

\Quote{那位天使向我宣告——在他的喜讯中，当他向我显现时——祂的国度没有终结——这奥秘被保守，不被揭示。}\par

\Quote{那星也向我们再次宣告——你儿子就是那位将持守冠冕者。——祂的容貌有所改变——祂本是天使，却未将\Emph{它}让我们知晓。}\par

\Quote{在我面前，当守望者显现时——祂在成孕之前就称祂为主——向我宣告祂是至高者之子——但祂的父在哪里，祂并未让我知道。}\par

\Quote{在我们面前，祂以星的形体宣告——那至高者的主就是所生的那位——你的儿子是光之星的主宰——除非祂命令，它们不升起。}\par

\Quote{看，在你面前——又启示了别的奥秘，好让你认识真理——我如何在童贞中生了我的儿子——祂是天主子，去，传扬祂！}\par

\Quote{在我们面前，那星教导我们——祂的诞生超越世界，你的儿子超越万有——祂是天主子，正如你所说的。}\par

\Quote{高天与世界之下都为祂作证——所有\OriginalTerm{守望者}{Watchers}和众星——祂是天主子，是主。——将祂的名声传到你们的土地！}\par

\Quote{整个高天，藉一颗星——激动了波斯，她认识了真理——你的儿子是天主子——万民都当臣服于祂。}\par

\Quote{将平安带到你们的土地——愿平安在你们的疆界倍增！——作为真理的宗徒，愿你们被相信——在你们所经过的一切路上。}\par

你圣子的平安，它必护送我们——安然归回我们的土地，正如它引领我们\Emph{前来}——当祂的能力握住世界时——愿祂眷顾我们的土地，赐福它！\par

\Quote{愿波斯因你的喜讯而喜乐！——愿亚述因你的来临而欢欣——当我儿子的国度兴起时——愿祂的旗帜插在你们的土地上！}\par

让教会欢欣歌唱——\Quote{荣耀归于至高者的诞生，——上界与下界都因祂而光照！}——愿祂受赞美，因祂的诞生使万有喜乐！\par

\SourceRecord{Translated by A. Edward Johnston.；Nicene and Post-Nicene Fathers, Second Series；Vol. 13.；Edited by Philip Schaff and Henry Wace.；Buffalo, NY: Christian Literature Publishing Co.,；1898.；<http://www.newadvent.org/fathers/3704.htm>.}
